\documentclass[11pt,a4paper]{ctexart}

\usepackage[margin=1in]{geometry}
\usepackage{amsmath,amssymb,booktabs,array,longtable}
\usepackage{hyperref}
\usepackage{xcolor}
\usepackage{enumitem}
\usepackage{listings}

\hypersetup{
  colorlinks=true,
  linkcolor=blue!55!black,
  urlcolor=blue!55!black,
  citecolor=blue!55!black
}

\lstset{
  basicstyle=\ttfamily\small,
  breaklines=true,
  frame=single,
  columns=fullflexible
}

\title{GNNHAR-IV 资产规模扩展实验方案：从 Dow 30 到 S\&P 100 / S\&P 500}
\author{Volatility Project Design Note}
\date{\today}

\begin{document}
\maketitle

\section{核心研究问题}

你现在这篇文章的主问题可以自然扩展成一个更强的 scale-effect 问题：

\[
\text{当资产数量 } N \text{ 从 }30\text{ 增加到 }100\text{ 或 }500\text{ 时，图结构带来的预测增益是否变大？}
\]

更具体地说，不只问 GNNHAR-IV 是否比 HAR-IV 好，而是估计下面这些增益如何随 \(N\) 变化：
\[
\Delta^{GHAR}_{N}
= L(HAR_N)-L(GHAR_N),
\]
\[
\Delta^{GNN}_{N}
= L(GHAR_N)-L(GNNHAR_N),
\]
\[
\Delta^{GNN\text{-}IV}_{N}
= L(GHAR\text{-}IV_N)-L(GNNHAR\text{-}IV_N),
\]
其中 \(L(\cdot)\) 可以分别取 test MSE 或 test QLIKE。论文中最重要的表述不是“某一个模型赢了”，而是
\[
\Delta^{GNN\text{-}IV}_{100} > \Delta^{GNN\text{-}IV}_{30}
\quad\text{或}\quad
\Delta^{GNN\text{-}IV}_{500} > \Delta^{GNN\text{-}IV}_{100}.
\]
这才直接回答“更大的图是否让图神经网络更有用”。

\section{当前项目的数据口径}

当前 Dow 30 pipeline 使用三张宽表：

\begin{longtable}{lll}
\toprule
文件 & 形状 & 含义 \\
\midrule
\endhead
\texttt{merged\_rv\_data\_filled.csv} & \(\text{Date}\times N\) & 目标变量 \(v_{i,t}\)，项目里的 realized volatility / historical volatility 口径 \\
\texttt{merged\_iv\_data\_filled.csv} & \(\text{Date}\times N\) & 隐含波动率 \(x_{i,t}\)，当前使用 AlphaQuery 30-day IV Mean \\
\texttt{dow30\_daily\_returns\_2021\_2026.csv} & \(\text{Date}\times N\) & 日收益率 \(r_{i,t}\)，用于 GLASSO 图结构 \\
\bottomrule
\end{longtable}

代码当前采用重叠 HAR 窗口：
\[
H^{RV}_{i,t-1}
=
\left(
v_{i,t-1},
\frac{1}{5}\sum_{\ell=1}^{5} v_{i,t-\ell},
\frac{1}{22}\sum_{\ell=1}^{22} v_{i,t-\ell}
\right),
\]
\[
H^{IV}_{i,t-1}
=
\left(
x_{i,t-1},
\frac{1}{5}\sum_{\ell=1}^{5} x_{i,t-\ell},
\frac{1}{22}\sum_{\ell=1}^{22} x_{i,t-\ell}
\right).
\]

扩展到 S\&P 100 或 S\&P 500 时，最重要的是保持同一口径：

\begin{itemize}[leftmargin=*]
\item \(v_{i,t}\) 不要一部分来自 LOBSTER 5-minute realized variance，另一部分来自 AlphaQuery historical volatility。
\item \(x_{i,t}\) 统一用 AlphaQuery 30-day IV Mean，必要时再加 call IV、put IV、skew 等扩展实验。
\item \(r_{i,t}\) 用同一调整价格源计算，例如 Yahoo Finance adjusted close 的 log return。
\item 时间范围固定，例如 2021-01-26 到 2026-01-26，所有 universe 取共同交易日。
\end{itemize}

\section{Universe 设计}

建议跑四个 universe，而不是只跑 30 和 500：

\begin{longtable}{lll}
\toprule
Universe & 目的 & 选择方法 \\
\midrule
\endhead
Dow 30 & 当前基准 & 保留现有结果，作为 \(N=30\) 的 anchor \\
S\&P 100 & 对齐 Zhang et al. 的大致规模 & 用 S\&P 500 中市值最大且 IV/RV 完整的 100 只 \\
S\&P 200 & 中间规模检验 & 用同一排序取前 200 只，观察增益曲线是否单调 \\
S\&P 500-clean & 最大规模检验 & 使用数据完整性过滤后的 S\&P 500 成分股 \\
\bottomrule
\end{longtable}

不建议直接说“全部 S\&P 500”而不筛选，因为 AlphaQuery IV 很可能存在缺失、换 ticker、并购退市、指数成分历史变化等问题。论文里应定义为：
\[
\mathcal{U}_{500}^{clean}
=
\{i\in \text{S\&P 500 current constituents}: \text{RV, IV, return coverage pass filters}\}.
\]

推荐过滤规则：

\begin{itemize}[leftmargin=*]
\item 时间覆盖率：每只股票在样本期内 \(v_{i,t}\)、\(x_{i,t}\)、\(r_{i,t}\) 的非缺失比例至少 \(95\%\)。
\item IV 极端值：winsorize 到每只股票 1\% 和 99\% 分位，或全局合理区间。
\item 股票数量固定：先筛出 clean universe，再按市值排序取 \(N=30,100,200,500\)。
\item Dow 30 作为单独结果保留；S\&P 30 可以作为 robustness，因为 Dow 30 与 S\&P top 30 不是同一组股票。
\end{itemize}

\section{数据来源方案}

\subsection{收益率}

收益率最容易：
\[
r_{i,t}=\log(P_{i,t}^{adj})-\log(P_{i,t-1}^{adj}).
\]
可用 Yahoo Finance 或 Stooq 获取 adjusted close。收益率只用于 GLASSO，不直接作为目标。

\subsection{IV}

沿用你当前 paper 的 AlphaQuery 口径：
\[
x_{i,t}= \text{AlphaQuery 30-day Implied Volatility Mean}.
\]

实际采集时每个 ticker 保存一份原始文件，再合并成宽表：

\begin{lstlisting}
data/raw/alphaquery_iv/AAPL.csv
data/raw/alphaquery_iv/MSFT.csv
...
data/processed/sp500/merged_iv_data_filled.csv
\end{lstlisting}

每个 ticker 的原始表至少保留：

\begin{lstlisting}
Date,ticker,iv_mean_30,iv_call_30,iv_put_30,
put_call_iv_ratio_30,iv_mean_skew_30,
historical_volatility_30
\end{lstlisting}

主实验只用 \texttt{iv\_mean\_30}。其他字段留给 appendix。

\subsection{RV / 目标变量}

这里有两个可选路线，必须二选一。

\paragraph{路线 A：和当前项目最一致。}
继续使用 AlphaQuery 或同源数据里的 historical/realized volatility 字段作为 \(v_{i,t}\)。优点是能覆盖 S\&P 500，和你当前 Dow 30 文件最接近。缺点是严格意义上不是自己从高频价格计算的 realized variance。

\paragraph{路线 B：更严格但更贵。}
从高频 5-minute returns 构造 realized variance：
\[
RV_{i,t}=\sum_{m=1}^{M} r_{i,t,m}^{2}.
\]
优点是更贴近 Zhang et al. 的 realized-volatility literature。缺点是 S\&P 500 高频数据成本和清洗工作量很大。

当前建议：主实验用路线 A，把变量名在论文里谨慎写成项目中的 volatility target \(v_{i,t}\)，不要过度声称是高频 realized variance；如果未来拿到高频数据，再做路线 B 的 robustness。

\section{模型组设计}

每个 universe 都跑同一套模型：

\begin{longtable}{lll}
\toprule
模型 & 输入 & 用途 \\
\midrule
\endhead
HAR & \(H^{RV}_{i,t-1}\) & 无图、无 IV 的基准 \\
GHAR & \(H^{RV}_{i,t-1}, (WH^{RV}_{t-1})_i\) & 线性图结构增益 \\
GNNHAR1L/2L/3L & \(H^{RV}_{t-1}, W\) & 非线性、多跳图增益 \\
HAR-IV & \(H^{RV}_{i,t-1}, H^{IV}_{i,t-1}\) & IV 信息基准 \\
GHAR-IV & \(H^{RV}, H^{IV}, WH^{RV}, WH^{IV}\) & 线性图和 IV 的互补性 \\
GNNHAR-IV & \([H^{RV},H^{IV}], W\) & 主要研究模型 \\
HAR-fakeIV & \(H^{RV}, \widetilde H^{IV}\) & IV 参数扩张 placebo \\
GHAR-fakeIV & \(H^{RV}, \widetilde H^{IV}, WH^{RV}, W\widetilde H^{IV}\) & 线性 placebo \\
GNNHAR-IV-fakeIV & \([H^{RV},\widetilde H^{IV}], W\) & 神经网络 placebo \\
GHAR-random / GNNHAR-random & 使用 density-matched random \(W\) & 图结构 placebo \\
\bottomrule
\end{longtable}

其中 fake IV 用日期置换：
\[
\widetilde H^{IV}_{i,t}=H^{IV}_{i,\pi(t)}.
\]
它保留 IV 的尺度和维度，但破坏时间预测关系。

\section{图结构}

训练期收益率矩阵记为 \(R_t\)。GLASSO 估计 precision matrix：
\[
\widehat \Theta
=
\arg\min_{\Theta\succ0}
\left\{
\operatorname{tr}(S\Theta)-\log\det(\Theta)
\lambda\|\Theta\|_1
\right\}.
\]

构造邻接矩阵：
\[
A_{ij}=\mathbf{1}\{|\widehat\Theta_{ij}|>\tau,\ i\ne j\}.
\]
归一化：
\[
W=D^{-1/2}AD^{-1/2}.
\]

S\&P 500 规模下 GLASSO 会更重，建议：

\begin{itemize}[leftmargin=*]
\item 先在 S\&P 100 和 S\&P 200 调好 \(\lambda\) 或 edge density。
\item S\&P 500 用固定目标密度，例如 \(5\%\)、\(10\%\)、\(15\%\)，避免图太密导致 \(WH\) 变成市场平均。
\item 主表报告选定图；appendix 报告不同 density 的 robustness。
\end{itemize}

\section{评估指标}

测试集 MSE：
\[
MSE
=
\frac{1}{|\mathcal{T}_{test}|N}
\sum_{t\in \mathcal{T}_{test}}\sum_{i=1}^{N}
(v_{i,t}-\widehat v_{i,t})^2.
\]

测试集 QLIKE：
\[
QLIKE
=
\frac{1}{|\mathcal{T}_{test}|N}
\sum_{t\in \mathcal{T}_{test}}\sum_{i=1}^{N}
\left(
\frac{v_{i,t}}{\widehat v_{i,t}}
-\log\frac{v_{i,t}}{\widehat v_{i,t}}-1
\right).
\]

核心结果要报告 loss ratio：
\[
\rho_{m,N}^{HAR}=\frac{L(m,N)}{L(HAR,N)},
\quad
\rho_{m,N}^{HAR\text{-}IV}=\frac{L(m,N)}{L(HAR\text{-}IV,N)}.
\]

再报告图结构增益：
\[
G^{linear}_N
=
1-\frac{L(GHAR_N)}{L(HAR_N)},
\]
\[
G^{gnn}_N
=
1-\frac{L(GNNHAR_N)}{L(GHAR_N)},
\]
\[
G^{gnn\text{-}iv}_N
=
1-\frac{L(GNNHAR\text{-}IV_N)}{L(GHAR\text{-}IV_N)}.
\]

如果 \(G^{gnn\text{-}iv}_{100}\) 或 \(G^{gnn\text{-}iv}_{500}\) 明显大于 \(G^{gnn\text{-}iv}_{30}\)，你的 scale hypothesis 才有强证据。

\section{代码组织}

建议新增以下目录，而不是继续把所有逻辑塞进 Dow 30 pipeline：

\begin{lstlisting}
data/
  raw/
    alphaquery_iv/
    prices/
  processed/
    dow30/
    sp100/
    sp200/
    sp500_clean/
configs/
  universes/
    dow30.yaml
    sp100.yaml
    sp200.yaml
    sp500_clean.yaml
scripts/
  data/
    build_universe.py
    fetch_prices.py
    fetch_alphaquery_iv.py
    build_panel.py
  analysis/
    gnnhar_iv_scale_pipeline.py
notebooks/
  gnnhar_iv_scale_colab.ipynb
outputs/
  scale_experiment/
    dow30/
    sp100/
    sp200/
    sp500_clean/
\end{lstlisting}

\subsection{配置文件}

每个 universe 一个 YAML：

\begin{lstlisting}
name: sp100
tickers_file: data/processed/sp100/tickers.txt
start_date: "2021-01-26"
end_date: "2026-01-26"
rv_file: data/processed/sp100/merged_rv_data_filled.csv
iv_file: data/processed/sp100/merged_iv_data_filled.csv
returns_file: data/processed/sp100/daily_returns.csv
split:
  train: 0.70
  valid: 0.15
  test: 0.15
graph:
  method: glasso
  target_density: 0.10
models:
  gnn_layers: [1, 2, 3]
  hidden_grid: [16, 32, 64]
  lr_grid: [0.0003, 0.001, 0.003]
\end{lstlisting}

\subsection{主运行接口}

主脚本接口应该变成：

\begin{lstlisting}
python scripts/analysis/gnnhar_iv_scale_pipeline.py \
  --config configs/universes/sp100.yaml \
  --output-dir /content/drive/MyDrive/GNNHAR-colab-runs/outputs/scale_experiment/sp100 \
  --tune-gnn \
  --mcs-bootstrap 300 \
  --seed 20260602
\end{lstlisting}

当前 \texttt{gnnhar\_iv\_pipeline.py} 已经有大部分模型逻辑。需要改的地方主要是：

\begin{itemize}[leftmargin=*]
\item \texttt{load\_panel()} 从固定 Dow 30 文件名改成读取 config。
\item \texttt{glasso\_adjacency()} 增加 target density 或固定 alpha 的选项。
\item \texttt{evaluate\_runs()} 增加 universe/name 字段，方便横向合并。
\item 新增 \texttt{scale\_summary.csv}，把不同 \(N\) 的 loss ratio 和 graph gain 汇总到一张表。
\end{itemize}

\section{Colab 运行方案}

Colab notebook 分成六个 cell 就够：

\begin{enumerate}[leftmargin=*]
\item mount Google Drive。
\item clone GitHub branch。
\item install requirements。
\item 检查每个 universe 的三张数据宽表。
\item 按 universe 循环运行 pipeline。
\item 合并结果，生成 scale-effect 表和图。
\end{enumerate}

核心命令：

\begin{lstlisting}
from google.colab import drive
drive.mount('/content/drive')

!git clone --depth 1 --branch 2026-06-01 https://github.com/easygl1der/GNNHAR.git /content/GNNHAR
%cd /content/GNNHAR

!pip install -q numpy pandas scikit-learn scipy matplotlib torch pyyaml

universes = ["dow30", "sp100", "sp200", "sp500_clean"]
for u in universes:
    !python scripts/analysis/gnnhar_iv_scale_pipeline.py \
      --config configs/universes/{u}.yaml \
      --output-dir /content/drive/MyDrive/GNNHAR-colab-runs/outputs/scale_experiment/{u} \
      --tune-gnn \
      --mcs-bootstrap 300 \
      --seed 20260602

!python scripts/analysis/summarize_scale_experiment.py \
  --input-root /content/drive/MyDrive/GNNHAR-colab-runs/outputs/scale_experiment \
  --output-dir /content/drive/MyDrive/GNNHAR-colab-runs/outputs/scale_experiment/summary
\end{lstlisting}

S\&P 500 的 GNNHAR3L-IV 可能比较慢。建议先跑 smoke test：

\begin{lstlisting}
python scripts/analysis/gnnhar_iv_scale_pipeline.py \
  --config configs/universes/sp100.yaml \
  --output-dir outputs/scale_experiment/sp100_smoke \
  --fast \
  --epochs 40
\end{lstlisting}

确认数据、维度、loss 正常后再跑 full run。

\section{最终论文表格}

主表建议按 universe 展开：

\begin{longtable}{lrrrrrr}
\toprule
Universe & \(N\) & HAR QLIKE & GHAR QLIKE & GNNHAR3L QLIKE & GHAR-IV QLIKE & GNNHAR3L-IV QLIKE \\
\midrule
\endhead
Dow 30 & 30 &  &  &  &  &  \\
S\&P 100 & 100 &  &  &  &  &  \\
S\&P 200 & 200 &  &  &  &  &  \\
S\&P 500-clean & \(N_{clean}\) &  &  &  &  &  \\
\bottomrule
\end{longtable}

第二张表直接报告增益：

\begin{longtable}{lrrrr}
\toprule
Universe & \(G^{linear}_N\) & \(G^{gnn}_N\) & \(G^{iv}_N\) & \(G^{gnn\text{-}iv}_N\) \\
\midrule
\endhead
Dow 30 &  &  &  &  \\
S\&P 100 &  &  &  &  \\
S\&P 200 &  &  &  &  \\
S\&P 500-clean &  &  &  &  \\
\bottomrule
\end{longtable}

其中
\[
G^{iv}_N=1-\frac{L(HAR\text{-}IV_N)}{L(HAR_N)}.
\]

\section{解释逻辑}

如果结果是：
\[
G^{gnn\text{-}iv}_{30}\leq 0
\quad\text{但}\quad
G^{gnn\text{-}iv}_{100}>0,
\]
你的解释会很强：Dow 30 过小，线性 GHAR 已经足够；更大图里，多跳和非线性 spillover 才有空间发挥。

如果结果是：
\[
G^{linear}_{N}>0
\quad\text{但}\quad
G^{gnn}_{N}\leq 0
\]
即使 \(N=500\)，那说明图结构有用，但 GNNHAR 不一定比线性 GHAR 稳定，论文就应该强调 GHAR-IV 是更稳健的工程方案。

如果 random graph 在大 \(N\) 下也接近 GLASSO，说明模型可能只是受益于 cross-sectional smoothing 或参数扩张，不应声称“金融网络被识别出来”。这时要依赖 GLASSO vs random 的 placebo 结果来保护论文解释。

\section{最小可执行路线}

建议分三阶段：

\begin{enumerate}[leftmargin=*]
\item 数据阶段：先完成 S\&P 100 的三张宽表，完全复用现有 pipeline 跑通。
\item 代码阶段：把 \texttt{gnnhar\_iv\_pipeline.py} 参数化成 config-driven pipeline。
\item 论文阶段：补一节 ``Network Scale Experiment''，把 Dow 30、S\&P 100、S\&P 200、S\&P 500-clean 的结果放进同一张表。
\end{enumerate}

最先跑 S\&P 100，因为它最接近 Zhang et al. 的原始规模，也最能回答你现在的问题；S\&P 500 是加分项，但数据完整性和 Colab 运行成本会更高。

\end{document}
